% --- Chapter: Multi-Dimensional Arrays ---
\section{Multi-Dimensional Arrays}

CRISP supports nested array literals, chained indexing, and nested iteration — all building on the existing array primitives without special syntax or additional types.

\subsection{Nested Array Literals}

Arrays can contain other arrays, enabling matrix-like and jagged structures:

\begin{tcolorbox}[colback=codebg, colframe=darkpurple, colbacktitle=darkpurple, title=\textbf{\textcolor{codegold}{Nested Arrays}}, fonttitle=\bfseries, sharp corners]
\begin{lstlisting}
# Single-line nested array
let @matrix = [[1, 2, 3], [4, 5, 6], [7, 8, 9]];

# Multi-line nested array
let @grid = [
    [1, 2, 3],
    [4, 5, 6],
    [7, 8, 9]
];

# Jagged arrays (rows of different lengths)
let @jagged = [[1, 2], [3, 4, 5], [6]];
\end{lstlisting}
\end{tcolorbox}

\subsection{Chained Indexing}

Postfix indexing chains naturally for multi-dimensional access:

\begin{tcolorbox}[colback=codebg, colframe=darkpurple, colbacktitle=darkpurple, title=\textbf{\textcolor{codegold}{Chained Indexing}}, fonttitle=\bfseries, sharp corners]
\begin{lstlisting}
let @matrix = [[1, 2, 3], [4, 5, 6], [7, 8, 9]];

# Single index returns a row (an array)
say @matrix[0];       # [1, 2, 3]

# Chained indexing drills into nested arrays
say @matrix[1][2];    # 6

# Deep nesting
let @cube = [[[1, 2], [3, 4]], [[5, 6], [7, 8]]];
say @cube[0][1][0];   # 3
\end{lstlisting}
\end{tcolorbox}

\clearpage
\subsection{Nested Iteration}

\texttt{for} loops iterate over each level, enabling nested traversal:

\begin{tcolorbox}[colback=codebg, colframe=darkpurple, colbacktitle=darkpurple, title=\textbf{\textcolor{codegold}{Nested For Loops}}, fonttitle=\bfseries, sharp corners]
\begin{lstlisting}
# Print each cell of a matrix
for row in @matrix {
    for cell in row {
        say cell;
    }
}
# Output: 1 2 3 4 5 6 7 8 9

# Collect all elements into a flat array
let @flat = [];
for row in @matrix {
    for cell in row {
        @flat.push(cell);
    }
}
say @flat;  # [1, 2, 3, 4, 5, 6, 7, 8, 9]
\end{lstlisting}
\end{tcolorbox}

\begin{tcolorbox}[colback=lightlavender, colframe=softvioletdark, title=\textbf{\textcolor{darkpurple}{Design Note: No Special Matrix Type}}, fonttitle=\bfseries, sharp corners]
CRISP does not have a special matrix or tensor type. Multi-dimensional arrays are simply arrays of arrays. This design choice provides:

\begin{itemize}
    \item \textbf{Simplicity:} No new types, syntax, or concepts to learn
    \item \textbf{Flexibility:} Jagged arrays (rows of different lengths) are natural
    \item \textbf{Composability:} All existing array methods (\texttt{.map()}, \texttt{.filter()}, \texttt{.contains()}) work on nested arrays without modification
\end{itemize}

Future releases may add dedicated linear algebra operations (\texttt{transpose}, \texttt{reshape}, \texttt{flatten}) as array methods.
\end{tcolorbox}
